\documentclass{article}
\usepackage{amsmath, amssymb, amsthm}
\usepackage{booktabs}
\usepackage{multirow}
\usepackage{geometry}
\usepackage{graphicx}
\usepackage{hyperref}
\usepackage{cleveref}
\usepackage{xcolor}

\geometry{margin=1in}
\hypersetup{colorlinks=true, urlcolor=blue}

% Shared macros for the multi-mode mixed-criticality scheduling papers.
% % Shared macros for the multi-mode mixed-criticality scheduling papers.
% % Shared macros for the multi-mode mixed-criticality scheduling papers.
% \input{macros} from any document in this directory rather than redefining
% these locally -- a macro defined in one paper's preamble and silently
% missing from a sibling's is exactly the class of error that put \round and
% \E undefined in evaluation_methodology.tex.

\newcommand{\N}{\mathbb{N}}
\newcommand{\R}{\mathbb{R}}
\newcommand{\Z}{\mathbb{Z}}
\newcommand{\E}{\mathbb{E}}
\newcommand{\Var}{\mathrm{Var}}
\newcommand{\Prob}{\mathbb{P}}
\newcommand{\Ind}{\mathbf{1}}
\newcommand{\ceil}[1]{\lceil #1 \rceil}
\newcommand{\floor}[1]{\lfloor #1 \rfloor}
\newcommand{\round}[1]{\lfloor #1 \rceil}
 from any document in this directory rather than redefining
% these locally -- a macro defined in one paper's preamble and silently
% missing from a sibling's is exactly the class of error that put \round and
% \E undefined in evaluation_methodology.tex.

\newcommand{\N}{\mathbb{N}}
\newcommand{\R}{\mathbb{R}}
\newcommand{\Z}{\mathbb{Z}}
\newcommand{\E}{\mathbb{E}}
\newcommand{\Var}{\mathrm{Var}}
\newcommand{\Prob}{\mathbb{P}}
\newcommand{\Ind}{\mathbf{1}}
\newcommand{\ceil}[1]{\lceil #1 \rceil}
\newcommand{\floor}[1]{\lfloor #1 \rfloor}
\newcommand{\round}[1]{\lfloor #1 \rceil}
 from any document in this directory rather than redefining
% these locally -- a macro defined in one paper's preamble and silently
% missing from a sibling's is exactly the class of error that put \round and
% \E undefined in evaluation_methodology.tex.

\newcommand{\N}{\mathbb{N}}
\newcommand{\R}{\mathbb{R}}
\newcommand{\Z}{\mathbb{Z}}
\newcommand{\E}{\mathbb{E}}
\newcommand{\Var}{\mathrm{Var}}
\newcommand{\Prob}{\mathbb{P}}
\newcommand{\Ind}{\mathbf{1}}
\newcommand{\ceil}[1]{\lceil #1 \rceil}
\newcommand{\floor}[1]{\lfloor #1 \rfloor}
\newcommand{\round}[1]{\lfloor #1 \rceil}

% Default parameter values quoted in evaluation_methodology.tex, defined once
% here rather than as literals scattered through the text.
%
% These MUST match amc_tasksim.generation.taskset.generate_taskset's actual
% keyword defaults -- tests/generation/test_paper_defaults.py asserts this by
% parsing the \newcommand definitions below and comparing them against
% inspect.signature(generate_taskset), so the paper cannot silently drift from
% what the code actually does. Update both together.

\newcommand{\defaultNumTasks}{20}
\newcommand{\defaultCP}{0.5}
\newcommand{\defaultU}{0.8}
\newcommand{\defaultInvFailProb}{10000}
\newcommand{\defaultTmin}{100}
\newcommand{\defaultTmax}{10000}
\newcommand{\defaultBCETmin}{0.8}


% --- Review-pass markup, added for a human review of this draft. Strip this
% block and every \reviewnote/\decisionnote call before final submission.
% Plain fcolorbox/parbox, not todonotes: todonotes' inline notes rendered as
% empty space in this document for reasons not worth chasing further, so this
% avoids the dependency rather than debugging it (see mode_optimization.tex's
% own copy of this note for the isolation trail). ---
\newcommand{\reviewnote}[1]{\par\noindent\fcolorbox{red!70!black}{red!8}{%
  \parbox{\dimexpr\linewidth-2\fboxsep-2\fboxrule\relax}{\textbf{REVIEW --} #1}}\par\smallskip}
\newcommand{\decisionnote}[1]{\par\noindent\fcolorbox{blue!60!black}{blue!6}{%
  \parbox{\dimexpr\linewidth-2\fboxsep-2\fboxrule\relax}{\textbf{DECISION --} #1}}\par\smallskip}

\title{Multi-Level Graceful Degradation for Mixed-Criticality Scheduling}
\author{}
\date{}

\begin{document}

\maketitle

\begin{center}
\fbox{\parbox{0.85\textwidth}{\footnotesize \textbf{Review copy, assembled from three
previously separate documents} (\texttt{task\_model.tex}, \texttt{evaluation\_methodology.tex},
\texttt{mode\_optimization.tex} -- each still compiles independently via the project
\texttt{Makefile}; this file is the combined reading order). Highlighted boxes flag specific
points for review (\textbf{\textcolor{red!70!black}{REVIEW}}) and open decisions for future work
(\textbf{\textcolor{blue!60!black}{DECISION}}). \Cref{sec:review-notes} near the end collects
items that do not attach to one location. Remove every \texttt{\textbackslash reviewnote}/%
\texttt{\textbackslash decisionnote} call (and their definitions, near the top of the source)
before any external submission.}}
\end{center}

\begin{abstract}
Mixed-criticality (MC) scheduling improves on worst-case pessimism by certifying
high-criticality (HI) tasks under an optimistic assumption during normal operation and falling
back to a pessimistic one if that assumption is violated. Across three decades of MC scheduling
research, that fallback has almost always been a single, abrupt step: some or all
low-criticality (LO) work is abandoned at once, and the literature's refinements have targeted
\emph{when} that step fires, not what happens once it does. Graceful degradation -- responding
in proportion to severity rather than as one cliff-edge -- is well established elsewhere in
dependable systems design, but has not been brought into MC scheduling with the same analytical
rigor the two-level case already has. This document presents a severity-indexed ladder of
degradation levels whose safety is certified by the same response-time analysis used for
classic two-level Adaptive MC (AMC) scheduling, generalised rather than replaced; the
evaluation methodology and metrics used to assess it; the staged empirical protocol used to
choose how the ladder should be configured, executed on real simulation data rather than merely
proposed; and, following directly from what that protocol found, a further result on \emph{exit}
timing -- when a degraded level should be left, not just entered -- that turned out to be the
largest single effect measured in this line of work. Results, methods, and a small number of
open decisions are reported together, on the position that a reviewer should be able to see
the whole argument in one document rather than reconstruct it from separate files.
\end{abstract}

\section{Introduction}
\label{sec:introduction}

Safety-critical systems increasingly host tasks of different criticality on the same platform --
functions whose failure could contribute to an accident alongside functions whose failure is
merely inconvenient. Certifying every task to the same, conservative worst-case execution time
(WCET) bound wastes most of the processor most of the time, since the bound a certification
authority will accept for a high-criticality (HI) function is typically far more pessimistic than
the execution time that function actually needs in normal operation. Vestal's original
formulation of mixed-criticality (MC) scheduling~\cite{Vestal} exploits this directly: certify
HI-criticality tasks against their pessimistic bound, but \emph{schedule} the system, in normal
operation, against a less pessimistic one -- and provide a fallback guarantee for the case where
a HI-criticality job actually needs more than that optimistic budget. The field this founded is
now large enough to have its own survey~\cite{MCS_survey} and its own recurring workshop, and the
question of what the fallback should look like, and how to justify it to a certification
authority, has been treated as a first-class problem in its own right rather than an
afterthought~\cite{graydon2013safetyassurance}.

\textbf{The fallback has almost always been a single, abrupt step.} The dominant formulation --
Adaptive MC (AMC), and the runtime co-design refinements this project's own prior work
contributes~\cite{9804600} -- has one normal mode and one degraded mode: a HI-criticality job
reaching a response-time threshold trips a single transition, after which low-criticality (LO)
work is abandoned, either entirely or via whatever conservative rule the exit protocol specifies.
Bailout protocols make the same shape more literal still, dropping every LO-criticality task at
the instant of detection rather than allowing any grading at
all~\cite{ecrts_bailout,tse_bailout}. Subsequent work in the field has continued to refine this
picture -- justifying more precisely how much LO-criticality service is actually
lost~\cite{rtns_2020_service}, and asking how robust the two-level guarantee is under realistic
assumptions~\cite{resilient_scheduling} -- but in every case the refinement targets \emph{when}
the abrupt step fires or how quickly the system returns from it, not what the step itself looks
like. Whether a HI-criticality overrun is barely over budget or catastrophically over budget, the
LO-criticality response today is the same: gone.

\textbf{Graceful degradation is not a new idea; it is simply new to this corner of scheduling
theory.} Responding to a fault or an overload in proportion to its severity, rather than as one
discrete fallback, has been a recognised attribute of dependable systems since Laprie's
foundational terminology for the field~\cite{Laprie_dependability_definition}, and a substantial,
independent line of work in fault-tolerant and distributed embedded systems has built and
evaluated staged degradation directly: fail-stop processors as a composable building block for a
graded response~\cite{schlichting83failstop}, general architectural frameworks for graceful
degradation in distributed embedded systems~\cite{shelton}, product-family approaches to the same
problem~\cite{nace}, and task-allocation algorithms explicitly optimised against a utility model
of how the system degrades as components fail~\cite{rtss_2008}. None of this is drawn from MC
scheduling -- it predates the field's current form by decades in some cases -- and none of it
carries a real-time response-time guarantee of the kind MC scheduling depends on. It establishes
something narrower but still useful here: that staged, proportionate degradation is a workable,
well-studied design pattern in systems that must keep working under fault, not a speculative one.

\textbf{Contribution.} This document brings graceful degradation into MC scheduling with the same
analytical grounding the two-level case already has, rather than as a heuristic utility function
layered on top of it. The mechanism is a severity-indexed ladder of intermediate degradation
levels, each protected by the \emph{same} response-time analysis already used to certify the
two-level guarantee -- generalised across levels, not replaced by a new argument per level
(\Cref{sec:design-rationale,sec:admissibility}). Three things follow from taking that route
seriously rather than assuming it works: a principled way to decide \emph{how many} levels and
where to place them, settled empirically rather than by convention
(\Cref{sec:optimisation-problem-statement}--\Cref{sec:protocol}); a full-scale empirical result
comparing the scheme against two-level AMC across utilisation and deadline regimes
(\Cref{sec:stage-results}); and a second, unanticipated result that the same machinery makes
answerable -- not just when to enter a degraded level, but how gracefully to \emph{leave} one,
which turns out to be the single largest effect measured anywhere in this document
(\Cref{sec:exit-results}).

\textbf{Structure of this document.} \Cref{sec:design-rationale}--\Cref{sec:notation-task-model}
present the task model, the safety argument for the graded ladder, and the two operating points
the model admits. \Cref{sec:task-set-generation}--\Cref{sec:notation-eval} present the evaluation
methodology: how task sets are generated, the metrics and objective function used to score a
configuration, and the simulation parameters held fixed throughout. From
\Cref{sec:optimisation-problem-statement} onward the document turns to the optimisation question
this line of work exists to answer -- how the ladder should be configured, and, following from
what that investigation found, how it should be exited -- reporting the staged protocol,
its results, and what they do and do not establish. This is a review draft: specific points are
flagged inline as they arise, and \Cref{sec:review-notes} near the end collects the small number
of items that concern the document as a whole.

\input{report_parts/task_model_body}

\input{report_parts/evaluation_methodology_body}

\input{report_parts/mode_optimization_body}

\section{Conclusion}
\label{sec:conclusion}

The central claim of this document is that graceful degradation can be added to mixed-criticality
scheduling without weakening the guarantee the field already relies on: every level of the
severity ladder is certified by the same response-time analysis used for the classic two-level
case, generalised rather than replaced (\Cref{sec:design-rationale,sec:admissibility}), and the
scheme beats two-level AMC-RA across every configuration of a properly-powered regime map
(\Cref{sec:stage-results}). Getting there required correcting more than one claim along the way
rather than only confirming them -- an underpowered pilot's apparent loss at the hardest tested
regime, a retracted unconditional safety statement replaced by a graded one, an overclaim about
whole-run time-degraded caught while writing its own proof -- and each correction is left visible
in the text rather than silently absorbed, on the view that a document which never revises itself
in front of the reader is less trustworthy, not more.

What was not anticipated at the outset is where the largest remaining effect turned out to live:
not in how many levels to use, or where to place their triggers -- both vacuous under the adopted
policy -- but in \emph{exit} timing, a question the original protocol was not designed to ask
(\Cref{sec:exit-results}). That result is reported with the same discipline as the rest of the
document: a real, measured cost (increased oscillation) alongside the real, measured benefit,
neither one asserted to dominate the other without a validated way to weigh them.

This is a review draft, not a finished one. \Cref{sec:review-notes} collects what still needs a
decision -- most of it now concentrated in exit-timing behaviour, not in the parts of the
argument that were settled earlier and have stood since.

\bibliographystyle{plain}
\bibliography{project_bib,references}

\end{document}
